\chapter{Background}
\label{chap:background}

This chapter introduces the three building blocks the thesis combines: the
Darwin G{\"o}del Machine (Section~\ref{sec:dgm}), multi-fidelity evaluation
schedulers (Section~\ref{sec:multifidelity}), and the SWE-bench evaluation
suite (Section~\ref{sec:swebench}).

\section{The Darwin G{\"o}del Machine}
\label{sec:dgm}

The Darwin G{\"o}del Machine~\cite{Sakana:25} is an evolutionary
self-improvement loop over coding agents.  An \emph{agent} is a directory
containing the Python source of an LLM-driven program that, given a problem
statement and a Git repository, attempts to produce a patch that solves the
problem.  Each generation proceeds in three phases:
%
\begin{enumerate}
  \item \textbf{Parent selection.}  A scoring rule picks parents from the
    archive of previously-kept agents.
  \item \textbf{Child generation.}  Each parent is asked to propose a
    modification to its own source code; the modified code becomes a
    \emph{child}.
  \item \textbf{Child evaluation.}  Each child is run against the benchmark.
    Children that pass quality and compile checks are added to the archive.
\end{enumerate}
%
The archive seeds the next generation, so improvements compound across
generations.

% TODO: cite original DGM paper formally, expand on archive policy and target
% selection (`score_child_prop`, contextlength targeting, etc.)

\section{Multi-Fidelity Evaluation}
\label{sec:multifidelity}

Multi-fidelity methods exploit the observation that cheap, partial
evaluations are often informative enough to discard the worst candidates
before paying for the full evaluation.  Successive halving runs all
candidates at a small budget, keeps the top $1/\eta$ fraction, and re-runs
them at a larger budget; the process repeats until one budget level
remains.  Hyperband~\cite{Li:18} composes successive halving over several
bracket sizes to balance exploration and exploitation.  ASHA~\cite{Li:20} is
an asynchronous variant in which a candidate is promoted to the next rung as
soon as it ranks in the top $1/\eta$ of the candidates that have completed
the current rung, eliminating the synchronous straggler problem.

In this thesis the candidates are children of a DGM generation and the
fidelities are the number of benchmark tasks each child is evaluated on.

\section{SWE-bench and \texttt{swe\_verified\_mini}}
\label{sec:swebench}

SWE-bench~\cite{Jimenez:24} is a benchmark of real GitHub issues from
popular Python projects.  Each task provides a problem statement, a base
commit, and a hidden set of tests that the agent's patch must make pass.
SWE-bench~Verified is a human-validated 500-task subset; we use a 50-task
sub-subset, \texttt{swe\_verified\_mini}, which is small enough to run
multiple full generations within the thesis budget while preserving the
distribution of task difficulties.

% TODO: describe per-task evaluation pipeline (Docker, conda envs, report.py)
